% GENERATED FILE. Edit canonical lessons, then run npm run generate:books.
% canonical-source-hash: fnv1a64:fc40da694426c1eb
% canonical-lessons: RU-C459-pokhod, RU-C459-tropa, RU-C459-prival, RU-C459-karabkatsya, RU-C459-spotykatsya

\chapter{Going hiking}
\label{ch:ru-going-hiking}
\hlchaptermodality{459}

\begin{chapteropening}
\textbf{By the end of this chapter:} \emph{I can go hiking: a hike, a trail, a rest stop, scrambling up and stumbling.}

Complete the last lesson of chapter 459: I can go hiking: a hike, a trail, a rest stop, scrambling up and stumbling.
\end{chapteropening}

\section[\texorpdfstring{\ru{поход} (pokhód)}{поход (pokhód)}]{\ru{поход} (pokhód) — a hike, a trek; a campaign}

\label{lesson:RU-C459-pokhod}



\begin{quote}
Before the new one: say the Russian for a summit, then the Russian for a gorge.
\end{quote}

\subsection*{You'll want to know: \ru{поход}}
\textbf{\ru{поход}} (\emph{pokhód}) — \textquotedblleft{}a hike, a trek; a campaign\textquotedblright{}.

\begin{quote}
\ru{Летом} \ru{мы} \ru{ходили} \ru{в} \ru{поход} \ru{в} \ru{горы}.
\end{quote}

\begin{quote}
\emph{Létom my khodíli v pokhód v góry.}
\end{quote}

\begin{quote}
In summer we went on a hike in the mountains.
\end{quote}

Masculine. \emph{\ru{Идти} \ru{в} \ru{поход}} — to go hiking or camping.

\begin{cousinweb}
\emph{\ru{поход}} is built from \emph{\ru{по}-} and \emph{\ru{ход}} 'a going', which is inherited from Proto-Slavic \emph{xodŭ} 'motion' (EtymDB).
\end{cousinweb}

\begin{grammarlens}[title={What you've built}]
One more word for telling, explaining and describing at B1.
\end{grammarlens}

\subsection*{Your turn}
\begin{itemize}
\raggedright
  \item \emph{Say it:} \emph{pokhód}
  \item \emph{Say it:} \emph{pokhód}, once more
  \item \emph{Say it:} a sentence of your own with \emph{pokhód}
  \item \emph{Recall:} say the Russian for a summit, then the Russian for a gorge, then say \emph{pokhód} again
\end{itemize}

\begin{tcolorbox}[breakable,colback=teal!4,colframe=teal!35!black,title={Before you move on}]
What does \ru{поход} mean? (\textquotedblleft{}A hike, a trek; a campaign\textquotedblright{}.) Say it once more.
\end{tcolorbox}
\section[\texorpdfstring{\ru{тропа} (tropá)}{тропа (tropá)}]{\ru{тропа} (tropá) — a path, a trail}

\label{lesson:RU-C459-tropa}



\begin{quote}
Before the new one: say the Russian for a gorge, then the Russian for a hike.
\end{quote}

\subsection*{You'll want to know: \ru{тропа}}
\textbf{\ru{тропа}} (\emph{tropá}) — \textquotedblleft{}a path, a trail\textquotedblright{}.

\begin{quote}
\ru{Узкая} \ru{тропа} \ru{ведёт} \ru{к} \ru{озеру}.
\end{quote}

\begin{quote}
\emph{Úzkaya tropá vedyót k ózeru.}
\end{quote}

\begin{quote}
A narrow path leads to the lake.
\end{quote}

Feminine; plural \emph{\ru{тропы}}. \emph{\ru{Тропинка}} — a little path.

\begin{cousinweb}
\emph{\ru{тропа}} is an old Slavic word, usually linked with a root for stamping the feet: a path trodden by feet. EtymDB lists it without a source, so the line is hedged.
\end{cousinweb}

\begin{grammarlens}[title={What you've built}]
One more word for telling, explaining and describing at B1.
\end{grammarlens}

\subsection*{Your turn}
\begin{itemize}
\raggedright
  \item \emph{Say it:} \emph{tropá}
  \item \emph{Say it:} \emph{tropá}, once more
  \item \emph{Say it:} a sentence of your own with \emph{tropá}
  \item \emph{Recall:} say the Russian for a gorge, then the Russian for a hike, then say \emph{tropá} again
\end{itemize}

\begin{tcolorbox}[breakable,colback=teal!4,colframe=teal!35!black,title={Before you move on}]
What does \ru{тропа} mean? (\textquotedblleft{}A path, a trail\textquotedblright{}.) Say it once more.
\end{tcolorbox}
\section[\texorpdfstring{\ru{привал} (privál)}{привал (privál)}]{\ru{привал} (privál) — a rest stop (on a hike), a halt}

\label{lesson:RU-C459-prival}



\begin{quote}
Before the new one: say the Russian for a hike, then the Russian for a path.
\end{quote}

\subsection*{You'll want to know: \ru{привал}}
\textbf{\ru{привал}} (\emph{privál}) — \textquotedblleft{}a rest stop (on a hike), a halt\textquotedblright{}.

\begin{quote}
\ru{Через} \ru{два} \ru{часа} \ru{мы} \ru{сделали} \ru{привал}.
\end{quote}

\begin{quote}
\emph{Chérez dva chasá my sdélali privál.}
\end{quote}

\begin{quote}
After two hours we stopped for a rest.
\end{quote}

Masculine. \emph{\ru{Сделать} \ru{привал}} — to stop for a rest.

\begin{cousinweb}
\emph{\ru{привал}} is built on \emph{\ru{привалить}} 'to lean against; to pull in to shore', from \emph{\ru{при}-} and \emph{\ru{валить}} 'to tumble'. EtymDB does not list it, so the line is hedged.
\end{cousinweb}

\begin{grammarlens}[title={What you've built}]
One more word for telling, explaining and describing at B1.
\end{grammarlens}

\subsection*{Your turn}
\begin{itemize}
\raggedright
  \item \emph{Say it:} \emph{privál}
  \item \emph{Say it:} \emph{privál}, once more
  \item \emph{Say it:} a sentence of your own with \emph{privál}
  \item \emph{Recall:} say the Russian for a hike, then the Russian for a path, then say \emph{privál} again
\end{itemize}

\begin{tcolorbox}[breakable,colback=teal!4,colframe=teal!35!black,title={Before you move on}]
What does \ru{привал} mean? (\textquotedblleft{}A rest stop (on a hike), a halt\textquotedblright{}.) Say it once more.
\end{tcolorbox}
\section[\texorpdfstring{\ru{карабкаться} (karábkatsya)}{карабкаться (karábkatsya)}]{\ru{карабкаться} (karábkatsya) — to scramble (up), to clamber}

\label{lesson:RU-C459-karabkatsya}



\begin{quote}
Before the new one: say the Russian for a path, then the Russian for a rest stop.
\end{quote}

\subsection*{You'll want to know: \ru{карабкаться}}
\textbf{\ru{карабкаться}} (\emph{karábkatsya}) — \textquotedblleft{}to scramble (up), to clamber\textquotedblright{}.

\begin{quote}
\ru{Дети} \ru{карабкаются} \ru{по} \ru{камням}.
\end{quote}

\begin{quote}
\emph{Déti karábkayutsya po kamnyám.}
\end{quote}

\begin{quote}
The children are scrambling over the rocks.
\end{quote}

Imperfective; perfective \emph{\ru{вскарабкаться}}. Present \emph{\ru{карабкаюсь}}.

\begin{cousinweb}
The origin of \emph{\ru{карабкаться}} is not settled; EtymDB lists the word without a source.
\end{cousinweb}

\begin{grammarlens}[title={What you've built}]
One more word for telling, explaining and describing at B1.
\end{grammarlens}

\subsection*{Your turn}
\begin{itemize}
\raggedright
  \item \emph{Say it:} \emph{karábkatsya}
  \item \emph{Say it:} \emph{karábkatsya}, once more
  \item \emph{Say it:} a sentence of your own with \emph{karábkatsya}
  \item \emph{Recall:} say the Russian for a path, then the Russian for a rest stop, then say \emph{karábkatsya} again
\end{itemize}

\begin{tcolorbox}[breakable,colback=teal!4,colframe=teal!35!black,title={Before you move on}]
What does \ru{карабкаться} mean? (\textquotedblleft{}To scramble (up), to clamber\textquotedblright{}.) Say it once more.
\end{tcolorbox}
\section[\texorpdfstring{\ru{спотыкаться} (spotykátsya)}{спотыкаться (spotykátsya)}]{\ru{спотыкаться} (spotykátsya) — to stumble, to trip}

\label{lesson:RU-C459-spotykatsya}



\begin{quote}
Before the new one: say the Russian for a rest stop, then the Russian for to scramble.
\end{quote}

\subsection*{You'll want to know: \ru{спотыкаться}}
\textbf{\ru{спотыкаться}} (\emph{spotykátsya}) — \textquotedblleft{}to stumble, to trip\textquotedblright{}.

\begin{quote}
\ru{В} \ru{темноте} \ru{я} \ru{часто} \ru{спотыкаюсь} \ru{о} \ru{камни}.
\end{quote}

\begin{quote}
\emph{V temnoté ya chásto spotykáyus o kámni.}
\end{quote}

\begin{quote}
In the dark I often trip over stones.
\end{quote}

Imperfective; perfective \emph{\ru{споткнуться}}. Takes \emph{\ru{о}} + accusative.

\begin{cousinweb}
\emph{\ru{спотыкаться}} is built from \emph{\ru{с}-}, \emph{\ru{по}-}, \emph{\ru{тыкать}} 'to poke, to knock' and \emph{-\ru{ся}} (EtymDB): to knock one's foot against something.
\end{cousinweb}

\begin{grammarlens}[title={What you've built}]
One more word for telling, explaining and describing at B1.
\end{grammarlens}

\subsection*{Your turn}
\begin{itemize}
\raggedright
  \item \emph{Say it:} \emph{spotykátsya}
  \item \emph{Say it:} \emph{spotykátsya}, once more
  \item \emph{Say it:} a sentence of your own with \emph{spotykátsya}
  \item \emph{Recall:} say the Russian for a rest stop, then the Russian for to scramble, then say \emph{spotykátsya} again
\end{itemize}

\begin{tcolorbox}[breakable,colback=teal!4,colframe=teal!35!black,title={Before you move on}]
What does \ru{спотыкаться} mean? (\textquotedblleft{}To stumble, to trip\textquotedblright{}.) Say it once more.
\end{tcolorbox}
